\subsection{компактность и центрированная систем замкнутных множетсв}
\begin{definition}[Компактное пространство]
	Метрическое пространство $X$ называется \textit{компактным}, если для любой системы открытых множеств $\{G_\alpha\}_{\alpha \in A} \subset 2^X$ такой, что $\bigcup\limits_{\alpha \in A} G_\alpha = X$, существует конечный набор $\alpha_1, \dotsc, \alpha_n \in A$ такой, что $\bigcup\limits_{k = 1}^n G_{\alpha_k} = X$.
\end{definition}

\begin{remark}
	Определение выше можно сформулировать так: из любого открытого покрытия пространства $X$ можно выделить конечное подпокрытие. Или, эквивалентно, если система открытых множеств не содержит конечного покрытия пространства $X$, то она не является покрытием пространства $X$.
\end{remark}

\begin{example}
	Любое замкнутое и ограниченное множество в $\mathbb{R}^n$ будет компактным (в индуцированной метрике).
\end{example}

\begin{definition}[Центрированная система]\label{center-system}
	Пусть $X$ --- метрическое пространство. Система множеств $\{B_\alpha\}_{\alpha \in A} \subset 2^X$ называется \textit{центрированной}, если для любого конечного набора индексов $\alpha_1, \dotsc, \alpha_n \in A$ выполнено:
	\[ \bigcap\limits_{k = 1}^n B_{\alpha_k} \neq \varnothing \]
\end{definition}



\begin{theorem}[Критерий компактности в терминах замкнутых множеств]\label{thm3.1}
	Следующие условия эквивалентны:
	\begin{itemize}
		\item Метрическое пространство $X$ компактно;
		\item Любая центрированная система замкнутых множеств в $X$ имеет непустое пересечение ($\bigcap_{\alpha \in A} B_\alpha \neq \varnothing$).
	\end{itemize}
\end{theorem}